\begin{resumo}

\noindent A análise dinâmica é uma etapa fundamental na engenharia para prever respostas vibratórias, garantir o conforto humano e prevenir falhas. O Método dos Elementos Finitos (MEF) tem sido uma das ferramentas mais utilizadas para esse propósito. Contudo, sua formulação clássica exige o uso de elementos com geometria restrita e com alta sensibilidade à distorção geométrica. O Método dos Elementos Virtuais (MEV) surge como uma generalização do MEF, capaz de lidar com elementos poligonais de topologia arbitrária, incluindo geometrias não convexas e nós colineares. O presente trabalho tem como objetivo investigar e implementar a formulação do MEV para a solução do problema de autovalores generalizados associado à análise de vibração livre de estruturas em estado plano. A principal vantagem da técnica reside na sua capacidade de operar sobre malhas não estruturadas formadas por polígonos genéricos, permitindo transições de malha naturais e adaptação facilitada aos contornos da estrutura. O arcabouço teórico adota as equações governantes da elastodinâmica, solucionadas numericamente com um programa desenvolvido especificamente para este fim. A viabilidade e a precisão do MEV foram validadas mediante a análise de três tipologias estruturais de interesse prático: uma chapa trapezoidal perfurada, a seção transversal de uma barragem de terra e o modelo de uma parede estrutural multi pavimentos. O seu desempenho foi medido por meio de métricas de erro em comparação aos resultados providos pelo MEF e pelo Método dos Elementos Finitos Generalizado (MEFG).

\end{resumo}
